\apendice{Documentación de usuario}

\section{Requisitos software y hardware para ejecutar el proyecto}
Para garantizar el correcto funcionamiento del sistema, se deben cumplir los requisitos mínimos indicados a continuación:

\subsection{Software}
\begin{itemize}
	\item \textbf{Sistema operativo:} Multiplataforma (Linux, Windows o macOS).
	\item \textbf{Lenguaje de programación:} Python 3.9 o superior.
	\item \textbf{Gestor de base de datos:} PostgreSQL 13 o superior (debido al uso de \texttt{psycopg2}).
	\item \textbf{Navegador web:} Versiones actualizadas de Google Chrome, Mozilla Firefox o Microsoft Edge (compatibles con HTML5 y CSS3).
\end{itemize}

\subsection{Hardware}
\begin{itemize}
    \item \textbf{Procesador:} Arquitectura x64.
    \item \textbf{Memoria RAM:} Mínimo 2 GB de RAM libres.
    \item \textbf{Almacenamiento:} 500 MB de espacio disponible para la aplicación y la base de datos local.
    \item \textbf{Conexión a internet:} Necesaria para el acceso remoto a la plataforma.
\end{itemize}

\section{Instalación / Puesta en marcha}
El proceso de instalación se basa en la creación de un entorno virtual para asegurar la estabilidad de las dependencias.

\subsection{Pasos para la ejecución en local}
\begin{enumerate}
	\item \textbf{Descarga del código:} Clonar el repositorio o descomprimir el archivo del proyecto.
	\item \textbf{Creación del entorno virtual:}
	\begin{verbatim}
	python -m venv venv
	venv\ Scripts\ activate
	\end{verbatim}
	\item \textbf{Instalación de dependencias:} Utilizar el archivo \texttt{requirements.txt} proporcionado:
	\begin{verbatim}
	pip install -r requirements.txt
	\end{verbatim}
	\item \textbf{Configuración de variables de entorno:} Crear un archivo \texttt{.env} en la raíz con las credenciales de la base de datos y la \texttt{SECRET\_KEY} (gestionado por \texttt{python-dotenv}).
	\item \textbf{Inicialización de la base de datos:} Asegurarse de que PostgreSQL está ejecutándose y crear la base de datos definida en el archivo de configuración.
	\item \textbf{Ejecución:}
	\begin{verbatim}
	flask run  # Desarrollo
	gunicorn -w 4 app:app  # Producción (usando gunicorn v26.0.0) \cite{gunicorn:manual}
	\end{verbatim}
\end{enumerate}

\section{Manuales y/o demostraciones prácticas}
La interacción con el sistema es intuitiva y sigue un flujo lineal dividido en tres áreas principales:

\subsection{Generación de recomendación}
El usuario accede a la página principal donde se presenta el cuestionario clínico. 
\begin{itemize}
	\item Debe completar todos los campos obligatorios sobre capacidades (motoras, cognitivas, visuales).
	\item Es imprescindible marcar la casilla de aceptación de términos legales para que el sistema procese la solicitud. \cite{ue:rgpd}
	\item Al finalizar, el sistema mostrará una página con los sistemas SAAC recomendados y el hardware complementario.
\end{itemize}

\subsection{Envío de retroalimentación (Feedback)}
Tras visualizar los resultados, el usuario encontrará un botón de Enviar Feedback. Al pulsarlo:
\begin{itemize}
	\item Se le realizarán preguntas sobre la exactitud de los resultados de recomendación.
	\item Podrá valorar la precisión de la recomendación de 1 a 10.
	\item Podrá añadir comentarios adicionales sobre su experiencia actual con otros sistemas.
	\item Esta información se vincula de forma anónima al registro previo mediante el identificador de sesión.
\end{itemize}

\subsection{Panel de administración (Uso del investigador)}
Accediendo a la ruta protegida \texttt{/admin\_historial}, el personal autorizado puede visualizar una tabla con todos los registros:
\begin{itemize}
	\item Se diferencian visualmente los "Pacientes reales" de las "Pruebas".
	\item Se pueden analizar los datos JSON almacenados para estudios clínicos o mejora del algoritmo de filtrado.
\end{itemize}
    
     